\section{Introduction}

% 1. Impact of Large lanugage model: 
    % Revolutionize human-like intelligence
% 2. LLM-based agent: 
    % Why LLM-based agent; to maximize the intelligence, LLM-based agent is proposed and developed.
    % What can LLM-based agent do / The importance of LLM-based Agent

% Impact of Large lanugage model: 
Large Language Models (LLMs) have recently shown remarkable potential in reaching a level of reasoning and planning capabilities comparable to humans.
% Why LLM-based agent;
This ability exactly aligns with the expectations of humans for autonomous agents that can perceive the surroundings, make decisions, and take actions in response~\cite{xi2023rise,Wooldridge1995IntelligentAT,Artificial_Intelligence_AModernApproach,guo2023indeed,liang2023let}. Hence, LLM-based agent has been studied and rapidly developed to understand and generate human-like instructions, facilitating sophisticated interactions and decision-making in a wide range of contexts
%to maximize this potential of LLMs, LLM-based agent has been studied and rapidly developed. 
% What is LLM-based agent
%The LLM-based agent refers to an autonomous agent powered by the LLM, enabling it to understand and generate human-like instructions, facilitating sophisticated interactions and decision-making in a wide range of contexts.
% LLM-based agent: previous work
%The LLM-based agent topic has attracted a lot of interests and various works have been developed by researchers
~\cite{yao2023tree,shinn2023reflexion,li2023apibank}. Timely survey papers systematically summarize the progress of LLM-based agents, as seen in works ~\cite{xi2023rise,wang2023survey}.
%Several systematic survey works~\cite{xi2023rise,wang2023survey} were also conducted to summarize the LLM-based agent. 


% \begin{figure}[tp]
% \begin{flushleft} % This will align the subgroup and caption to the left
% LLM-based Multi-Agents
% \begin{subgroup}
%   Idea Principal \\[2em]
%   Idea Principal \\[2em]
%   Idea Principal
% \end{subgroup}
% \caption{Caption describing the relationship between the general idea and the main ideas.}
% \end{flushleft}
% \end{figure}

\begin{figure*}[ht]
\begin{center}
\includegraphics[scale=0.56]{img/Trend_example.png}
\end{center}
\caption{The rising trend in the research field of LLM-based Multi-Agents. For Problem Solving and World Simulation, we categorize current work into several categories and count the number of papers of different types at 3-month intervals. The number at each leaf node denotes the count of papers within that category.}
\label{trend}
\end{figure*}


% 3. LLM-based Multi-Agent:
    % Why we need Multi-Agent
    % What is Multi-Agent
    % Why we focus on the Multi-Agent (importance of multi-agent)

% Why we need LLM-based Multi-Agent: Collective Intelligence
Based on the inspiring capabilities of the single LLM-based agent, LLM-based Multi-Agents have been proposed to leverage the collective intelligence and specialized profiles and skills of multiple agents. 
% What is LLM-based Multi-Agent 
\xz{Compared to systems using a single LLM-powered agent, multi-agent systems offer advanced capabilities by \underline{1)} specializing LLMs into various distinct agents, each with different capabilities, and \underline{2)} enabling interactions among these diverse agents to simulate complex real-world environments effectively. In this context, multiple autonomous agents collaboratively engage in planning, discussions, and decision-making, mirroring the cooperative nature of human group work in problem-solving tasks. This approach capitalizes on the communicative capabilities of LLMs, leveraging their ability to generate text for communication and respond to textual inputs. Furthermore, it exploits LLMs' extensive knowledge across various domains and their latent potential to specialize in specific tasks. Recent research has demonstrated promising results in utilizing LLM-based multi-agents for solving various tasks,}
%The LLM-based Multi-Agents refers to a system where multiple autonomous agents are empowered by LLM. These agents utilize the capabilities of LLM to interact with each other or the environment of the system. 
% Why we focus on the Multi-Agent (importance of multi-agent, research trend: not AI but interdisciplinary researchers => This Multi-Agent have broader impact to many fields)
%Compared to the individual agent, in the LLM-based Multi-Agents system, each agent's specialized capabilities contribute to the system's ability to efficiently scale up and adapt to changing environments 
such as software development~\cite{hong2023metagpt,qian2023communicative}, multi-robot systems~\cite{mandi2023roco,zhang2023building}, %etc. Additionally, the communication between agents in the Multi-Agents system enables the simulation of complex scenarios such as 
society simulation~\cite{park2023generative,park2022social}, policy simulation~\cite{xiao2023simulating,hua2023war}, and game simulation~\cite{xu2023language,wang2023avalon}. 
\xz{Due to the nature of interdisciplinary study in this field, it has attracted a diverse range of researchers, expanding beyond AI experts to include those from social science, psychology, and policy research. The volume of research papers is rapidly increasing,}
%Attracted by the powerful capabilities of the LLM-based Multi-agents systems, an increasing variety of researchers, extending beyond just AI researchers to those in fields such as social science, psychology, and policy researchers have delved into this promising research direction 
as shown in Fig. \ref{trend} (inspired by the design in ~\cite{gao2023retrieval}),
%\footnote{\scriptsize The way our figure is designed is also inspired by ~\cite{gao2023retrieval}.}, 
thus broadening the impact of LLM-based Multi-Agent research.
% Limitations (No summary)
Nonetheless, earlier efforts were undertaken independently,  resulting in an absence of a systematic review to summarize them, establish comprehensive \xz{blueprint of this field}, and examine future research challenges. This underscores the significance of our work and serves as the motivation behind presenting this survey paper, dedicated to the research on LLM-based multi-agent systems.


% 4. Objective of this survey:
    % Our purpose and contribution: AI researchers, Interdisciplinary researchers
      % how our survey aligns with current trends in AI and LLM research.

We expect that our survey can make significant contributions to both the research and development of LLMs and to a wider range of interdisciplinary studies employing LLMs. Readers will gain a comprehensive overview of  LLM-based Multi-Agent (LLM-MA) systems, grasp the fundamental concepts involved in establishing multi-agent systems based on LLMs, and catch the latest research trends and applications in this dynamic field. We recognize that this field is in its early stages and is rapidly evolving with fresh methodologies and applications. To provide a sustainable resource complementing our survey paper, we maintain an open-source GitHub repository\footnote{\scriptsize \url{https://github.com/taichengguo/LLM_MultiAgents_Survey_Papers}}.
We hope that our survey will inspire further exploration and innovation in this field,
%LLM-based Multi-Agent techniques and 
as well as applications across a wide array of research disciplines.

%Given the tremendous effort in this rapidly developing research area, we are motivated to present a comprehensive and systematic survey focusing on the LLM-based Multi-Agents. 
% Our purpose and contribution: AI researchers, Interdisciplinary researchers
%We expect that our survey can contribute to both LLM research fields and the broader interdisciplinary studies of LLM-based Multi-Agents, encompassing areas such as society, psychology, and policy, etc. For researchers who are already familiar with LLM-based Multi-Agents concepts,  our survey presents the latest applications to demonstrate current research trends. For newcomers or Non-CS researchers who are interested in the use of LLM-based Multi-Agents for problem solving or world simulation, we provide systematic taxonomy of LLM-based Multi-Agents to help them obtain a quick grasp of the core concepts of LLM-based Multi-Agents.
% Potential future Challenges and developments
%We are the view that this field is still in its early phase and we hope that our survey will motivate further exploration and innovation in LLM-based Multi-Agent techniques and applications across a diverse spectrum of research disciplines.

% 5. The detail structure of the survey paper
% The specific aspects of our survey: 
% How we organize the current work: Scope / Taxonomy, ......     Potential future Challenges and developments
%We organize this survey in the following structure. Related works are presented in Section~\ref{Sec2:background}. By drawing connections among the existing LLM-based Multi-Agents research, we first provide a detailed taxonomy, which summarizes four key properties (Agents-Environment Interface, Agents Profiling, Agents Communication and Agents Capabilities Acquisition) of the LLM-based Multi-Agents systems in Section~\ref{Sec3:Taxonomy}. Then we summarize the LLM-based Multi-Agents Framework in Section~\ref{Sec4:Framework} and Multi-Agents Orchestration and Efficiency in Section~\ref{Sec5:Orchestration}. We categorize current applications into two main streams: Multi-Agents for problem solving and Multi-Agents for World Simulation in Section~\ref{Sec5:Application}. We then look into the current LLM-based Multi-Agents datasets and benchmarks in Section~\ref{Sec6:Datasets and Benchmarks}. Based on the previous summary, we open the discussion for future research challenges and opportunities in Section~\ref{Sec7:Challenges and Opportunities}. The conclusions are summarized in Section~\ref{Sec8:Conclusion}.

To assist individuals from various backgrounds in understanding LLM-MA techniques and to complement existing surveys by tackling unresolved questions, we have organized our survey paper in the following manner. After laying out the background knowledge in Section~\ref{Sec2:background}, we address a pivotal question: \emph{How are LLM-MA systems aligned with the collaborative task-solving environment}?   To answer this, we present a comprehensive schema for positioning, differentiating, and connecting various aspects of LLM-MA systems in  \textbf{Section~\ref{Sec3:Taxonomy}}. We delve into this question by discussing: 1) the  \textbf{agents-environment interface}, which details how agents interact with the task environment; 2)  \textbf{agent profiling}, which explains how an agent is characterized by an LLM to behave in specific ways; 3)  \textbf{agent communication}, which examines how agents exchange messages and collaborate; and 4)  \textbf{agent capability acquisition}, which explores how agents develop their abilities to effectively solve problems. 
An additional perspective for reviewing studies about LLM-MA is their application. In  \textbf{Section~\ref{Sec5:Application}}, we categorize current  \textbf{applications} into two primary streams: multi-agents for \textbf{problem-solving} and multi-agents for \textbf{world simulation}.
To guide individuals in identifying appropriate \textbf{tools and resources},  we present open-source implementation frameworks for studying LLM-MA, as well as the usable datasets and benchmarks in  \textbf{Section~\ref{Sec4:Framework}}.
Based on the previous summary, we open the discussion for future research challenges and opportunities in Section~\ref{Sec7:Challenges and Opportunities}. The conclusions are summarized in Section~\ref{Sec8:Conclusion}.